\section{Quantum phase constraints}
\label{subsec:x9-mirror-periods}

Mirror symmetry supplies closed-M2 indices independently of the
magnetic sheaf counts. Let $n^0_\beta$ be the genus-zero
Gopakumar--Vafa invariant of the resolved threefold in curve class
$\beta$. We fix its normalization by
\begin{equation}
 \mathcal F_{\rm inst}(t)
 =-\frac{1}{(2\pi i)^3}\sum_{\beta>0}n^0_\beta
       \operatorname{Li}_3(e^{2\pi i t\cdot\beta}).
 \label{eq:x9-gv-prepotential}
\end{equation}
These are integer coefficients of a spin-weighted M2 index, not
dimensions of magnetic sheaf moduli
\cite{Gopakumar:1998BPS}. Applying the toric Frobenius construction
\cite{Hosono:1995Mirror,Demirtas:2024Mirror} gives
\begin{equation}
 n^0_{C_1}=n^0_{C_2}=1,\qquad n^0_R=0,\qquad
 n^0_{\iota_*e_1}=n^0_{\iota_*e_2}=1,
 \label{eq:x9-electric-gv}
\end{equation}
where $\iota:E\hookrightarrow\wideX_9$ and
$e_1=D_4|_E$, $e_2=D_3|_E$. The period construction and integral marking are given in
\cite[Section~5 and Appendix~C]{Wuebben:2026quantum}.
It determines all coefficients through total nef degree eight:
of $3002$ nonzero lattice classes in the chosen outer Mori cone,
$80$ have nonzero genus-zero invariant.

\paragraph{The mixed electric class has no connected curve.}
The vanishing in \eqref{eq:x9-electric-gv} also has a geometric
proof, valid beyond the finite degree computation. Let
$\pi:\wideX_9\to X_9$ be the contraction, whose exceptional fibers
are the disjoint loci $C_1$, $C_2$ and $E$. For any $m,n>0$, a
connected stable map in class $\beta=mC_1+nC_2$ would have zero
degree against the pullback of an ample divisor on $X_9$.
Its image must therefore lie in one fiber of $\pi$.
It cannot lie in either nodal fiber, since $C_1,C_2$ are independent
classes and both coefficients are positive. It cannot lie in $E$:
\begin{equation}
 E\cdot\beta=0,\qquad E|_E=K_E,
 \qquad -K_E\ \text{is ample}.
 \label{eq:x9-mixed-map-obstruction}
\end{equation}
Every nonconstant effective curve in $E$ has strictly negative
intersection with $E$. Hence the connected stable-map space is
empty for this class in every genus. In particular
$n^0_{mC_1+nC_2}=0$ for all $m,n>0$. This argument concerns
connected maps; the disconnected cycle $C_1\sqcup C_2$ certainly
exists.

The other entries of \eqref{eq:x9-electric-gv} are the isolated
$(-1,-1)$ rational-curve contributions, reproduced by the mirror
calculation. The compact relation now has a spectral interpretation:
\begin{equation}
 \iota_*e_1=\iota_*e_2+C_1+C_2,
 \qquad
 J\cdot\iota_*e_1=J\cdot\iota_*e_2+J\cdot C_1+J\cdot C_2.
 \label{eq:x9-m2-threshold-relation}
\end{equation}
An isolated rational curve in $E$ has the same total compact charge
and five-dimensional BPS mass as the disconnected three-curve
configuration on the right. The genus-zero single-particle index
at the left-hand charge is $+1$, whereas the difference class $R$
has no connected holomorphic representative. This is a threshold
charge and mass identity, not a calculation of a decay amplitude.
It explains why the $+1$ magnetic contribution at relative charge
$(R,-1)$ cannot be read as an elementary particle of charge $R$.
Neither the vanishing nor this threshold identity determines the
interacting operator spectrum at the singular point.

\paragraph{Computed magnetic corrections.}
Let $H_1,\ldots,H_6$ be the integral nef basis specified in Section~\ref{sec:probe-setup}, and let $\beta_a$ be its dual curve basis.
Thus $\beta_1=C_2$, $\beta_5=C_1$, $\beta_6=\iota_*e_2$, and
$\iota_*e_1=\beta_1+\beta_5+\beta_6$.
Write $t=\sum_a t_aH_a$ and $q_a=e^{2\pi i t_a}$.
Differentiating \eqref{eq:x9-gv-prepotential} gives the corrections
in \eqref{eq:x9-probe-quantum-period} in the same Mukai frame:
\begin{equation}
 \Psi_P(t)=\frac{1}{4\pi^2}\sum_{\beta>0}
      (P\cdot\beta)n^0_\beta\operatorname{Li}_2(q^\beta).
 \label{eq:x9-computed-magnetic-period}
\end{equation}
No additional polynomial or $c_2$ term is included here. Since
$A=-H_1+H_2$ and $V=H_1-H_2-H_5+H_6$, the first terms are
\begin{equation}
 \begin{aligned}
 4\pi^2\Psi_A&=-q_1+q_2+\tfrac14(-q_1^2+q_2^2)
                     +3q_2q_3-2q_2q_6+O(q^3),\\
 4\pi^2\Psi_V&=q_1-q_2-q_5+q_6
          +\tfrac14(q_1^2-q_2^2-q_5^2+q_6^2)
                     -3q_2q_3+q_3q_6+O(q^3).
 \end{aligned}
 \label{eq:x9-magnetic-period-leading}
\end{equation}
The remainder denotes total degree at least three, not powers of a
single coordinate. The one-quarter coefficients are double covers
of the degree-one classes.

To evaluate the previous path, use the flat coordinates
\begin{equation}
 (t_1,\ldots,t_6)=
 (-s+i\lambda\epsilon,\ i\lambda\epsilon,\ i\lambda(b+\epsilon),\
 i\lambda(a+\epsilon),\ i\lambda\epsilon,\ i\lambda\epsilon).
 \label{eq:x9-mirror-flat-path}
\end{equation}
Let $s_{[d]}$ denote the root of the alignment equation after
retaining GV terms of total degree at most $d$ in
\eqref{eq:x9-computed-magnetic-period}, with the full dilogarithm
for each retained class. At $t=-sD_5+iJ_*$ the degree-eight values are
\begin{equation}
 \begin{array}{c@{\quad}r@{\quad}r@{\quad}r}
  &s^{\rm pol}&s_{[8]}&s_{[8]}-s^{\rm pol}\\ \hline
 \mathcal F_2&0.740626137605&0.740632681385&6.54378\times10^{-6}\\
 \mathcal F_R&0.793227613297&0.793233852071&6.23877\times10^{-6}
 \end{array}
 \label{eq:x9-quantum-wall-numbers}
\end{equation}
The degree-six and degree-eight results differ by less than
$2\times10^{-20}$ here. This is agreement between truncations,
not a rigorous bound on the omitted classes. The truncated roots
are transverse and satisfy the same-phase inequality.

These numbers compare finite truncations near large volume.
They are not used to infer the phase sign at the transition.

\subsection{Charge transport and circle holonomy}
There is an exact continuation statement near either ordinary
conifold, with the other sectors kept nonsingular. Put
$\tau_i=t\cdot C_i$ and choose the vanishing brane
$\mathcal O_{C_i}(-1)$, whose charge is $(0,0,C_i,0)$.
The isolated-curve contribution is
\begin{equation}
 \Psi_P^{(i)}=\frac{P\cdot C_i}{4\pi^2}
        \operatorname{Li}_2(e^{2\pi i\tau_i})
 =\frac{P\cdot C_i}{2\pi i}\tau_i\log(-2\pi i\tau_i)
       +\text{single-valued terms}.
 \label{eq:x9-nodal-period-monodromy}
\end{equation}
A counterclockwise loop around $\tau_i=0$ consequently adds
$(P\cdot C_i)\tau_i$ to $Z$. In the original period branch this
is equivalent to the integral charge transformation
\begin{equation}
 M_i:(P,q,q_0)\longmapsto
       (P,q+(P\cdot C_i)C_i,q_0).
 \label{eq:x9-nodal-charge-monodromy}
\end{equation}
We use this period-continuation convention; transforming the charge
basis instead gives the inverse map. This is the usual conifold
transvection, also described by the spherical-twist monodromy
\cite[Section~5]{Aspinwall:2017Discriminants}.

The two transformations commute, since the vanishing charges are
mutually local. They are not equivalent to a single conifold
transvection associated with $R$. On the integral magnetic--electric
lattice,
\begin{equation}
 \begin{aligned}
 (M_1M_2-1)(P,q)&=
       \bigl(0,(P\cdot C_1)C_1+(P\cdot C_2)C_2\bigr),\\
 (M_R-1)(P,q)&=\bigl(0,(P\cdot R)R\bigr),\\
 \operatorname{rank}(M_1M_2-1)&=2,\qquad
 \operatorname{rank}(M_R-1)=1.
 \end{aligned}
 \label{eq:x9-two-node-monodromy}
\end{equation}
Here $M_R$ is a hypothetical single-particle transvection, not an
asserted discriminant monodromy. This distinction is consistent
with $n^0_R=0$ and is not removed by the compact homology relation.
For the magnetic constituents the needed pairings are
\begin{equation}
 \begin{array}{c@{\qquad}rr}
 P&P\cdot C_1&P\cdot C_2\\\hline
 A&0&-1\\ V&-1&1\\ T&-1&0
 \end{array}.
 \label{eq:x9-nodal-probe-pairings}
\end{equation}
Thus $M_2$ changes the electric charges of $\mathcal G$ and
$\mathcal H$ by $-C_2$ and $+C_2$, respectively, while leaving
their total probe charge fixed. The combined nodal monodromy shifts
the probe and its seed by $-C_1$, rather than $-R$. Their relative
charges remain unchanged. These formulas determine charge transport,
not stability or an index at the common singular point.

On \eqref{eq:x9-mirror-flat-path}, the normalized electric periods
of the two nodal D0 towers are
\begin{equation}
 Z_{1,k}=k+i\lambda\epsilon,\qquad
 Z_{2,k}=k-s+i\lambda\epsilon,\qquad k\in\ZZ.
 \label{eq:x9-nodal-tower-endpoint}
\end{equation}
At fixed auxiliary-circle radius, simultaneous nodal masslessness
therefore requires $\epsilon\to0$ and $s\to\ZZ$.
If the alignment tends to an interior value $s_*\in(0,1)$ instead,
the second tower has a nonzero normalized mass gap
$\operatorname{dist}(s_*,\ZZ)$. Thus shrinking the classical curve
areas along this $B$-field path does not by itself reach the
simultaneous massless locus. This is a condition on the auxiliary
four-dimensional probe problem, not an obstruction to the original
five-dimensional geometric transition.


\subsection{The marked compact transition}
Write $(v,a,b)=(z_2,z_3,z_4)$ for the spectator parameters.
For sufficiently small nonzero values, the companion paper proves
that the locus
\begin{equation}
 \Sigma:\quad z_1=z_5=z_6=1
 \quad\Longleftrightarrow\quad t_1=t_5=t_6=0
 \label{eq:x9-finite-bulk-locus}
\end{equation}
has four ordinary nodes. In the continued large-volume marking,
their integral charges are precisely
\begin{equation}
 (0,0,C_1,0),\quad(0,0,C_2,0),\quad
 (0,0,\iota_*e_1,0),\quad(0,0,\iota_*e_2,0).
 \label{eq:x9-finite-bulk-integral-charges}
\end{equation}
They generate a primitive rank-three lattice, with the single
relation $\iota_*e_1-\iota_*e_2-C_1-C_2=0$.
This assertion uses convergent electric periods and unit local
residues, not just the abstract relation or Hodge numbers
\cite[Section~5.2 and Appendix~C]{Wuebben:2026quantum}.
Its analytic proof is not repeated here.

In this compact integral marking, the magnetic pairings with the
exceptional-curve doublet are
\begin{equation}
 \begin{array}{c@{\qquad}rr}
 P&P\cdot\iota_*e_1&P\cdot\iota_*e_2\\\hline
 A&-1&0\\V&1&1\\T&0&1
 \end{array}.
 \label{eq:x9-local-probe-pairings}
\end{equation}
Thus the monodromy around this pair shifts the electric charge of
$\mathcal G$ by $-\iota_*e_1$, that of $\mathcal H$ by
$\iota_*e_1+\iota_*e_2$, and the total probe and its seed each
by $\iota_*e_2$. Their relative charges remain unchanged.
With $\tau_i=Z_{\gamma_i}$, the corresponding singular part of
a magnetic period due to these two nodes is
\begin{equation}
 Z_P^{\rm sing}=
 \frac1{2\pi i}\sum_{i=1}^2(P\cdot\iota_*e_i)
                 \tau_i\log\tau_i.
 \label{eq:x9-local-magnetic-singular-part}
\end{equation}
The two $C_1,C_2$ nodes contribute the analogous terms with their
previously computed magnetic pairings. Each singular term tends
to zero at the massless point and does not determine the
regular, bulk-dependent part of either constituent period.
In particular it does not fix
$\operatorname{Im}(Z_{\mathcal G}\overline{Z_{\mathcal H}})$.
These singular pieces alone do not settle the phase alignment.
The regular terms can, however, be continued on the same
small-bulk branch.

\subsection{Regular periods and the constituent phase}
The same continuation gives the regular magnetic periods
\cite[Appendix~C]{Wuebben:2026quantum}. In the Mukai frame of
\eqref{eq:x9-probe-quantum-period}, their leading terms are
\begin{equation}
 \begin{aligned}
 \Psi_A|_\Sigma&=-\frac1{12}-\frac{v+a}{4\pi^2}
                          +O(\delta^2),\\
 \Psi_V|_\Sigma&=\frac1{12}+\frac{2(v+a)}{4\pi^2}
                          +O(\delta^2),
 \qquad \delta=|v|+|a|+|b|.
 \end{aligned}
 \label{eq:x9-regular-magnetic-leading}
\end{equation}
Here $\Sigma$ denotes \eqref{eq:x9-finite-bulk-locus}; the error
is a small-bulk analytic remainder, not a total nef-degree
truncation. The constant terms include the four conifold
multicovers. The convergence proof in the companion paper sums all local
degrees before expanding in the spectator parameters.

For positive sufficiently small $(v,a,b)$, all three spectator
flat coordinates are purely imaginary. Write
$t_j=iy_j$ for $j=2,3,4$, with $y_j>0$, and put
$J_\Sigma=y_2H_2+y_3H_3+y_4H_4$. The restricted magnetic
corrections are real. The two quotient charges differ by $C_1$
with zero D0 shift, so their periods coincide exactly on $\Sigma$.
Using their Mukai charges gives
\begin{equation}
 \begin{aligned}
 Z_{\mathcal G}&=\mathcal V_A+\tfrac12+\Psi_A-iY,\\
 Z_{\mathcal H_2}=Z_{\mathcal H_R}
               &=\mathcal V_V-\tfrac5{12}+\Psi_V+iY,
 &Y&=y_3+\tfrac32y_4,
 \end{aligned}
 \label{eq:x9-endpoint-constituent-periods}
\end{equation}
where $\mathcal V_P=PJ_\Sigma^2/2$. In particular, the common
total probe period is real and equals
\begin{equation}
 \begin{aligned}
 M_\Sigma&=\mathcal V_T+\tfrac1{12}+\Psi_T,
 &\Psi_T&=\Psi_A+\Psi_V,\\
 \mathcal V_T&=2y_2y_3+2y_3^2+3y_2y_4+8y_3y_4+7y_4^2.
 \end{aligned}
 \label{eq:x9-endpoint-total-period}
\end{equation}
It is strictly positive when the positive bulk parameters are
sufficiently small: $y_3,y_4\to+\infty$, whereas $\Psi_T\to0$.
For either specified constituent split the exact alignment
function therefore satisfies
\begin{equation}
 \operatorname{Im}(Z_{\mathcal G}\overline{Z_{\mathcal H}})
       =-Y M_\Sigma<0.
 \label{eq:x9-endpoint-phase-sign}
\end{equation}
Thus the target is not on this marginal-stability wall. With
$\langle\Gamma_{\mathcal G},\Gamma_{\mathcal H}\rangle=-1$,
the necessary two-center stability inequality is satisfied
\cite[eq.~(3.23)]{Denef:2007Halos}; equivalently the relative
Fayet--Iliopoulos parameter has the bound-state sign used in
\eqref{eq:x9-probe-quiver-wall}.
This conclusion is a phase test for the specified charges.
It does not prove survival of both constituent BPS states from
large volume, exclude other decay channels, or extend the
two-node quiver index to a locus with four additional massless
particles. Those questions remain necessary for a compact
magnetic bound-state index.

\subsection{Light-particle emission and confinement}
\label{subsec:x9-magnetic-higgs}

Fix a positive sufficiently small bulk point on $\Sigma$. Approach it
in flat coordinates with the spectator coordinates fixed and
\begin{equation}
 t_1=i\epsilon c_2,\qquad t_5=i\epsilon c_1,\qquad
 t_6=i\epsilon c_e,\qquad c_1,c_2,c_e>0,\quad\epsilon\downarrow0.
 \label{eq:x9-light-electric-ray}
\end{equation}
The four light periods, ordered as $(C_1,C_2,\iota_*e_1,\iota_*e_2)$,
are exactly $i\epsilon(c_1,c_2,c_1+c_2+c_e,c_e)$. At the endpoint,
\eqref{eq:x9-endpoint-constituent-periods} gives positive real parts
for $Z_{\mathcal G}$ and both $Z_{\mathcal H}$ when the bulk
parameters are sufficiently small. The same holds for both total
probe periods. Continuity, including the vanishing
$\tau\log\tau$ terms, preserves these inequalities on a final
segment $0<\epsilon<\epsilon_0$.

Let $\Gamma$ be any of these five charges and let
$\eta=\sum_i n_i\gamma_i$, with $n_i\geq0$ and not all zero,
be a sum of the four light charges. Then
$Z_\eta=i\epsilon m_\eta$, where
$m_\eta=n_1c_1+n_2c_2+n_3(c_1+c_2+c_e)+n_4c_e>0$, and
\begin{equation}
 \operatorname{Im}(Z_{\Gamma-\eta}\overline{Z_\eta})
 =-\epsilon m_\eta\operatorname{Re}Z_\Gamma<0.
 \label{eq:x9-light-emission-sign}
\end{equation}
Thus $\Gamma\to(\Gamma-\eta)+\eta$ cannot be marginal on this
segment. The CPT-conjugate electric charge has the opposite, also
nonzero, alignment. This excludes these light-emission walls, whether
or not the proposed residual charge supports a BPS state. It neither
excludes other magnetic decompositions nor supplies a stable state at
the beginning of the segment. At $\epsilon=0$ the electric masses
vanish, so this argument does not define an endpoint particle index.

\paragraph{Confined magnetic flux.}
Choose the orientations
\begin{equation}
 (\delta_1,\delta_2,\delta_3,\delta_4)
 =(-\gamma_{C_1},-\gamma_{C_2},\gamma_{e_1},-\gamma_{e_2}),
 \qquad \sum_i\delta_i=0.
 \label{eq:x9-oriented-condensate}
\end{equation}
At a generic Higgs vacuum all four hypermultiplets have nonzero
expectation values. A D4 magnetic charge $P\in H^2(\wideX_9,\ZZ)$
then has flux winding vector
$w(P)=(\langle P,\delta_i\rangle)_{i=1}^4$, where the pairing
is $P\cdot\gamma_i$ before the orientation change. Its image lies in
\begin{equation}
 \Lambda_{\rm flux}=\{w\in\ZZ^4:\textstyle\sum_iw_i=0\}.
 \label{eq:x9-confined-flux-lattice}
\end{equation}
This is the $A_3$ root lattice as an integral winding lattice; it
does not assert an $SU(4)$ gauge enhancement. Write
$P=\sum_a p^aH_a$ in the primitive nef basis. The independent pairings are
$(P\cdot C_1,P\cdot C_2,P\cdot\iota_*e_2)=(p^5,p^1,p^6)$.
They range independently over $\ZZ^3$, so the image is all of
\eqref{eq:x9-confined-flux-lattice}, with no finite index.
The matter charges generate a saturated lattice; consequently
the unbroken subgroup of the $U(1)^6$ gauge torus is connected.
The $\ZZ_4$ quotient of the Higgs-branch coordinates is not an
unbroken discrete gauge group at a generic vacuum.

The relevant magnetic pairings are
\begin{equation}
 \begin{array}{c|rrrr|rrrr}
 &P C_1&P C_2&P\iota_*e_1&P\iota_*e_2
 &w_1&w_2&w_3&w_4\\\hline
 A&0&-1&-1&0&0&1&-1&0\\
 V&-1&1&1&1&1&-1&1&-1\\
 T&-1&0&0&1&1&0&0&-1\\
 D_8&0&0&1&1&0&0&1&-1\\
 E&0&0&-1&-1&0&0&-1&1
 \end{array}.
 \label{eq:x9-full-probe-flux}
\end{equation}
Here $A=D_4$, $V=D_7+D_8$ and $T=D_6=A+V$, as above.
The fluxes of the two constituents do not cancel: their sum is
$w(T)=(1,0,0,-1)$. Thus neither constituent nor either total
probe can continue as an isolated particle with that magnetic
charge at a generic Higgs vacuum. Its broken magnetic flux must
be carried by strings. This is the usual confinement mechanism
in a conifold Higgs transition \cite[Sections~3--4]{Greene:1996Confinement}.
Adding any electric charge leaves this obstruction unchanged.
In the mirror surgery, the same pairings specify the boundary
of the three-chain obtained from a magnetic three-cycle.
Neither the existence of a particular BPS core before Higgsing
nor the tension or BPS character of its flux strings follows
from this charge calculation.
